% ============================================
% Johns Hopkins Letter Template
% Assistant Professor Cover Letter – Centre College
% ============================================

\documentclass[11pt,letterpaper]{letter}

\usepackage{graphicx}
\usepackage{eso-pic}
\usepackage{geometry}
\usepackage{microtype}
\usepackage[hidelinks]{hyperref}

% ----- Page Setup -----
\geometry{left=25mm,right=25mm,top=25mm,bottom=25mm}

% ----- Logo Placement (first page only) -----
\newcommand{\PutJHULogoFG}[1]{%
  \AddToShipoutPictureFG*{%
    \AtPageUpperLeft{%
      \hspace{10mm}%
      \raisebox{-38mm}{%
        \includegraphics[width=6cm]{#1}%
      }%
    }%
  }%
}

% ----- Sender Information -----
\signature{Decory Edwards}
\address{Department of Economics \\ Johns Hopkins University \\ Baltimore, MD 21218 \\ \texttt{dedwar65@jh.edu}}

\begin{document}
\PutJHULogoFG{Images/jhulogo.png}

\begin{letter}{Search Committee \\
Economics \& Business Program \\
Centre College \\
Danville, KY 40422 \\ USA}

\opening{Dear Members of the Search Committee,}

I am writing to express my interest in the tenure-track Assistant Professor position in the Economics \& Business Program at Centre College, beginning August 2026, with a specialization in macroeconomics. I am a Ph.D. candidate in Economics at Johns Hopkins University, advised by Professors Christopher D.~Carroll, Nicholas W.~Papageorge, and Stelios Fourakis, and I expect to receive my degree in May 2026. My research fields are macroeconomics, wealth and income dynamics, and computational economics.

My research agenda focuses on understanding the determinants of wealth inequality and the mechanisms through which heterogeneity in returns to wealth shapes aggregate outcomes. In my job-market paper, I estimate the degree of heterogeneity in individual portfolio returns using micro-data from the Survey of Consumer Finances and simulate a structural model in which households face idiosyncratic return risk. I show that allowing for persistent heterogeneity in returns substantially improves the model’s ability to match the empirical distribution of wealth, offering a tractable way to connect micro-level return dispersion to macroeconomic inequality.

In a second project, I use the Health and Retirement Study to examine the relationship between interpersonal trust and portfolio performance. Building on the literature that finds a hump-shaped relationship between trust and income, I show that a similar relationship holds for individual returns to wealth measured in the HRS data, highlighting a behavioral channel through which social attitudes shape saving and investment outcomes. This work also provides new estimates of the persistent component of returns to net worth, which motivate the calibration of heterogeneity in my structural model.

Centre College’s teacher-scholar model, commitment to celebrating differences as stated in the Statement on Community, and emphasis on inclusive and varied pedagogical practices align closely with my approach to undergraduate education. I am particularly excited by the opportunity to contribute to core macroeconomics courses and electives in my areas of specialization, including public finance, open-economy macro, and related fields. The Economics \& Business Program’s dual-degree structure, strong mentoring culture, and exceptionally engaged student body make Centre an ideal fit for my interests in macroeconomic teaching, research, and undergraduate advising. I also value the College’s dedication to global learning and student development—especially its nationally recognized study abroad programs—and its support for early-career faculty through mentoring and professional development.

\textbf{Teaching.} My teaching experience centers on undergraduate macroeconomics and an upper-level macro policy seminar, where I have led weekly discussion sections and held regular office hours for classes ranging from 16 to 40 students. My approach is student-centered: I aim to help students “convince themselves” that they have moved from confusion to understanding by holding structured office-hour coaching that builds trust and confidence. At Centre, I am prepared to teach \textit{Principles of Macroeconomics}, \textit{Intermediate Macroeconomic Theory}, \textit{Money and Banking}, and \textit{Econometrics}, along with electives in macroeconomic policy, public finance, or open-economy macroeconomics. I welcome the opportunity to work closely with students in small-class environments and to support their academic growth through research mentorship and individualized advising.

I believe I would be a strong fit because my computational and empirical toolkit allows me to (i) provide rigorous, data-driven instruction across macroeconomics and quantitative methods, (ii) integrate reproducible workflows and real-world economic applications into my courses, and (iii) mentor students effectively in a liberal-arts environment that values inclusivity, curiosity, and high-impact teaching. I am excited about the opportunity to contribute to Centre College’s teaching, research, and community-engaged missions.

I would be honored to join the Economics \& Business Program at Centre College. Thank you for your consideration; I welcome the opportunity to discuss my fit with your department at your convenience.

\closing{Sincerely,}

\end{letter}
\end{document}
